% Consolidated textbook-ch05 from a4681eacb57a1f9fed79c1da26e5c6d0418a7d1c; source fragments preserved.

% === BEGIN textbook-ch05-choices.tex blob a60401a723eca9acb0cc92f9e2f480fb251fe92a ===
% 《大学物理简明教程》(第5版)第5章“气体动理论”
% 选择题5.1.1—5.1.14，原书/扫描PDF第123—125页。
% 题意为研究性转述，保持原题模型；所有条目待独立二校。
\subsection{选择题：平衡态与分子运动}

\sourceproblem{5.1.1}{123}{123}
\problemtext 金属杆两端分别与恒温沸水、冰接触，杆内温度分布已经稳定，
但不同位置温度不同。判断金属杆是否处于热力学平衡态。
\workings
热力学平衡态要求宏观上没有持续热流且各位置温度相同。
稳态导热满足$\partial T/\partial t=0$，但温差导致傅里叶热流
\[
\bm q=-\kappa\nabla T\ne\bm 0.
\]
故它是\emph{非平衡定常态}，不是热力学平衡，原题选\textbf{B}。
\crosscheck 时间上温度不变并不意味着空间温度均匀；不能把“稳定”和“平衡”混用。

\sourceproblem{5.1.2}{123}{123}
\problemtext 分子质量$m$的理想气体处于温度$T$，求分子速度在$x$方向分量的均方值$\langle v_x^2\rangle$。
\workings
平衡态各平动方向分配相同平均能量，
$\frac12m\langle v_x^2\rangle=\frac12k_{\mathrm B}T$。
\[
\boxed{\langle v_x^2\rangle=\frac{k_{\mathrm B}T}{m}}\quad(\text{原题D}).
\]
\crosscheck 总均方速率$\langle v^2\rangle=3k_{\mathrm B}T/m$，
与一个方向的均方速度有因子3的差异。

\sourceproblem{5.1.3}{123}{123}
\problemtext 同一容器内有1 mol氧气与1 mol氢气，
两气体在同一体积$V$及温度$T$下比较各自分压。
\workings
理想气体分压定律：$p_iV=n_iRT$。因为$n_1=n_2=1$，
\[
\boxed{p_{\mathrm{O_2}}=p_{\mathrm{H_2}}=\frac{RT}{V}}
\]
（选\textbf{C}）。
\crosscheck 分压取决于粒子数而非分子质量；不同摩尔质量不影响同温同容同摩尔数的分压。

\sourceproblem{5.1.4}{124}{124}
\problemtext 判断关于气体温度的四种描述：温度表征平均平动动能、
是分子热运动的集体统计表现、反映热运动激烈程度、可以代表每一个分子的“冷热”。
\workings
由
\[
\left\langle \varepsilon_{\mathrm{trans}}\right\rangle
=\frac32k_{\mathrm B}T
\]
可知前三种宏观统计表述正确，但单个分子并无与宏观温度一一对应的确定动能。
所以正确为①②③，原题选\textbf{B}。
\crosscheck 同一温度下速率遵从分布，单个分子的速率、动能不断变化。

\sourceproblem{5.1.5}{124}{124}
\problemtext 温度、压强相同的氢气与氧气，比较分子平均\emph{总动能}
$\overline\varepsilon_k$和平均\emph{平动动能}$\overline w$（理想刚性双原子分子模型）。
\workings
两气体都是常温下可按刚性双原子分子处理，自由度$i=5$。
因此
\[
\overline\varepsilon_k=\frac52k_{\mathrm B}T,\qquad
\overline w=\frac32k_{\mathrm B}T
\]
对两种气体均相同，故\textbf{两者均相等（A）}。
\crosscheck 不能将氢气误当单原子气体；在常温刚性双原子近似下氢气和氧气具有相同自由度。

\sourceproblem{5.1.6}{124}{124}
\problemtext 1 mol刚性双原子分子理想气体温度为$T$，求其内能。
\workings
在转动自由度充分激发而振动自由度冻结的温度区域，
$i=3+2=5$，有
\[
\boxed{U=\frac{i}{2}RT=\frac52RT}\quad(\text{选C}).
\]
\crosscheck 这里是1 mol而非1个分子；若是每分子平均能量应使用$k_{\mathrm B}$而不是$R$。

\sourceproblem{5.1.7}{124}{124}
\problemtext 判断理想气体单位时间内分子对单位面积器壁施加的总冲量随
体积、压强、温度、密度变化的关系，选择必然成立者。
\workings
该冲量大小除以时间、面积即为压强$p$：
\[
\boxed{\frac{\mathrm dJ}{\mathrm dt\,\mathrm dS}=p}.
\]
因此压强增大时单位时间单位面积的总冲量一定增大，选\textbf{B}。
体积减小、温度不变、密度不变若缺少对其他状态变量的约束，
均不足以保证该冲量单调或保持不变。
\crosscheck 由$p=nk_{\mathrm B}T$可构造反例：
若$n$不变而$T$升高，压强仍可升高。

\sourceproblem{5.1.8}{124}{124}
\problemtext 一定体积内的理想气体温度升为原来的2倍，
比较分子平均平动动能与容器内压强。
\workings
温度以开尔文计，$\langle\varepsilon_{\mathrm{trans}}\rangle=\tfrac32 k_{\mathrm B}T$，
而$n=N/V$不变，$p=nk_{\mathrm B}T$。
故两者均变为原来的2倍，原题选\textbf{A}。
\crosscheck 若改为摄氏温度数值“加倍”则不成立；
原题中的温度倍数必须指绝对热力学温度。

\sourceproblem{5.1.9}{124}{124}
\problemtext 氢气、氧气温度相同，原题给出关于两气体压强、密度、
单个分子速率及方均根速率的四个“氧气一定更大”的判断。
\workings
仅知$T$相同而没有体积、总质量、分子数或压强条件，
不能推出哪种气体压强大或密度大。单个分子速度有分布，
不可能由温度保证每个氧分子速度都更大。
方均根速率实际上
\[
v_{\mathrm{rms}}=\sqrt{\frac{3RT}{M_{\mathrm{mol}}}},
\quad
\frac{v_{\mathrm{rms,O_2}}}{v_{\mathrm{rms,H_2}}}
=\sqrt{\frac{2}{32}}=\frac14<1.
\]
故按原题扫描中的A—D四个断言逐字解释，\textbf{均不能成立}。
\crosscheck 原题没有“都不正确”选项；保留\textbf{题面或选项待核实}标记，
不编造某个字母为正确答案。

\sourceproblem{5.1.10}{124}{124}
\problemtext 理想气体放在恒温容器中并缓慢漏气，判断压强、
每个分子平均平动动能和剩余气体总内能的变化。
\workings
恒温即$T$不变，所以
\[
\boxed{\langle\varepsilon_{\mathrm{trans}}\rangle
=\frac32k_{\mathrm B}T=\text{常量}}.
\]
但容器内分子数$N$减少，$p=Nk_{\mathrm B}T/V$减少，
剩余气体的总内能$U=(i/2)Nk_{\mathrm B}T$也减少。
因此只有平均平动动能不变的描述正确，选\textbf{B}。
\crosscheck 单分子的平均能量与整个容器的内能是不同物理量。

\sourceproblem{5.1.11}{124}{124}
\problemtext 经典平衡理想气体中，比较最概然速率$v_p$、
平均速率$\overline v$和方均根速率$v_{\mathrm{rms}}$。
\workings
麦克斯韦分布给出
\[
v_p=\sqrt{\frac{2k_{\mathrm B}T}{m}},\quad
\overline v=\sqrt{\frac{8k_{\mathrm B}T}{\pi m}},\quad
v_{\mathrm{rms}}=\sqrt{\frac{3k_{\mathrm B}T}{m}}.
\]
因$2<8/\pi<3$，
\[
\boxed{v_p<\overline v<v_{\mathrm{rms}}}\quad(\text{选C}).
\]
\crosscheck 三者均有速率量纲，且随温度升高按$\sqrt T$增长。

\sourceproblem{5.1.12}{124--125}{124--125}
\problemtext 麦克斯韦速率分布图$f(v)$在$v_0$处达到峰值，
图中峰两侧区域$A,B$的面积标为相等。判断$v_0$的物理意义。
\workings
曲线在$v_0$处取最大值，即$f'(v_0)=0$且为峰点，
故$v_0$是\textbf{最概然速率}，选\textbf{A}。
\[
\boxed{v_0=v_p=\sqrt{\frac{2k_{\mathrm B}T}{m}}}
\]
（经典三维麦克斯韦速率分布）。
\crosscheck 实际严格麦克斯韦速率密度通常并非以众数为中心左右面积相等；
图中“$A,B$面积相等”若为严格条件，与标准分布峰点有潜在不一致。
因此保留图示\textbf{待二校}，不将众数误认为中位数。

\sourceproblem{5.1.13}{125}{125}
\problemtext 速率概率密度$f(v)$的物理意义为何？
\workings
$f(v)\mathrm dv$表示分子速率处于$[v,v+\mathrm dv]$区间的概率或分子数比例，
而$f(v)$本身是\emph{单位速率区间的概率密度}，故选\textbf{B}。
\[
\boxed{\frac{\mathrm dP}{\mathrm dv}=f(v),\qquad
P(v<v_i<v+\mathrm dv)=f(v)\mathrm dv}.
\]
\crosscheck $[f(v)]=(\mathrm{m/s})^{-1}$，不是无量纲百分率，也不是分子数。

\sourceproblem{5.1.14}{125}{125}
\problemtext 一定量理想气体在温度固定时降低压强，判断
平均碰撞频率$\overline Z$与平均自由程$\overline\lambda$的变化。
\workings
硬球气体近似下
\[
\overline\lambda=\frac{k_{\mathrm B}T}
 {\sqrt2\pi d^2p},\qquad
\overline Z=\frac{\overline v}{\overline\lambda}.
\]
$T$不变使$\overline v$不变；压强$p$减小时
$\overline\lambda$增大，而$\overline Z$减小。
故\textbf{碰撞频率减小、自由程增大（D）}。
\crosscheck 压强与自由程反比仅适用经典稀薄硬球模型；
稠密气体的多体关联需额外修正。
% === END textbook-ch05-choices.tex ===

% === BEGIN textbook-ch05-blanks.tex blob 6f81a275f82081200d8977a58182e041e721db8b ===
% 《大学物理简明教程》第5章填空题5.2.1—5.2.8。
% 书内及用户扫描PDF第125页。非官方独立推导；待二校。
\subsection{填空题：气体状态、能量与速率分布}

\sourceproblem{5.2.1}{125}{125}
\problemtext 一只电子管的真空度（管内气体压强）为
$1.0\times10^{-5}\,\mathrm{mmHg}$，温度为$27^\circ\mathrm C$。
求单位体积的气体分子数。
\workings
将压强换为SI单位：
$p=(10^{-5})(133.322)\,\mathrm{Pa}=1.33322\times10^{-3}\,\mathrm{Pa}$，
$T=300.15\,\mathrm K$。
使用$p=nk_{\mathrm B}T$与$k_{\mathrm B}=1.380649\times10^{-23}\,\mathrm{J/K}$，
\[
\boxed{n=\frac p{k_{\mathrm B}T}
\approx3.22\times10^{17}\,\mathrm{m^{-3}}}.
\]
\crosscheck 单位$\mathrm{Pa}/(\mathrm{J/K}\cdot\mathrm K)=\mathrm{m^{-3}}$。

\sourceproblem{5.2.2}{125}{125}
\problemtext 容器体积$10^{-2}\,\mathrm{m^3}$，其中气体质量为$100$ g，
分子的方均根速率为$200\,\mathrm{m/s}$；求压强。
\workings
密度$\rho=M/V=0.1/0.01=10\,\mathrm{kg/m^3}$，
由分子动理论压强公式
\[
\boxed{p=\frac13\rho\langle v^2\rangle
=\frac13(10)(200)^2
\approx1.33\times10^5\,\mathrm{Pa}}.
\]
\crosscheck $\rho v^2$量纲为$\mathrm{kg/(m\,s^2)}=\mathrm{Pa}$。

\sourceproblem{5.2.3}{125}{125}
\problemtext 温度$127^\circ\mathrm C$、1 mol氧气，
氧分子按刚性双原子分子处理。求气体内能与分子转动总动能。
\workings
$T=400.15$ K，氧分子有3平动与2转动自由度，所以
\[
\boxed{U=\frac52RT
\approx8.32\times10^3\,\mathrm J},\qquad
\boxed{E_{\mathrm{rot}}=RT\approx3.33\times10^3\,\mathrm J}.
\]
\crosscheck 总内能应大于纯转动能量；振动自由度在本题的刚性分子模型中未计入。

\sourceproblem{5.2.4}{125}{125}
\problemtext 标准状态下1立方米任意理想气体所含分子数。
\workings
标准状态取$p_0=1.01325\times10^5\,\mathrm{Pa}$、
$T_0=273.15\,\mathrm K$，由$pV=Nk_{\mathrm B}T$，
\[
\boxed{N=\frac{p_0(1\,\mathrm{m^3})}{k_{\mathrm B}T_0}
\approx2.69\times10^{25}}.
\]
\crosscheck 阿伏加德罗常量与摩尔体积之比
$N_A/(22.4\times10^{-3}\,\mathrm{m^3})$给近似同一数量级。

\sourceproblem{5.2.5}{125}{125}
\problemtext 三种容器分别装1 mol氮气、1 mol氢气与1 mol氦气，
绝对温度都升高1 K；求三种气体内能增量（刚性理想气体）。
\workings
氮气与氢气为双原子，$i=5$；氦气为单原子，$i=3$。
固定摩尔数时$\Delta U=(i/2)nR\Delta T$，故
\[
\boxed{\Delta U_{\mathrm{N_2}}=\frac52R
\approx20.79\,\mathrm J},\quad
\boxed{\Delta U_{\mathrm{H_2}}=\frac52R
\approx20.79\,\mathrm J},
\]
\[
\boxed{\Delta U_{\mathrm{He}}=\frac32R\approx12.47\,\mathrm J}.
\]
\crosscheck 刚性双原子比单原子每摩尔多两个转动自由度，故增量比为$5:3$。

\sourceproblem{5.2.6}{125}{125}
\problemtext 温度300 K时1 mol氢气分子的总平均平动动能、总平均转动动能与气体内能。
\workings
在刚性双原子近似下
\[
\boxed{E_{\mathrm{trans}}=\frac32RT\approx3.74\times10^3\,\mathrm J},\quad
\boxed{E_{\mathrm{rot}}=RT\approx2.49\times10^3\,\mathrm J},
\]
\[
\boxed{U=\frac52RT\approx6.24\times10^3\,\mathrm J}.
\]
\crosscheck 两类总动能之和等于气体内能；这里“总平均”是将分子的平均能量对1 mol求和。

\sourceproblem{5.2.7}{125}{125}
\problemtext 图5.2.7给出两条麦克斯韦速率分布：(1)峰值高、靠左、窄，
(2)峰值低、靠右、宽。判断同一种气体何者温度高，
以及同温氢气与氧气何者对应哪条曲线。
\workings
麦克斯韦分布速率尺度满足$v\propto\sqrt{T/m}$。
对同一种气体，温度高时曲线右移并展宽，故
\[
\boxed{\text{温度较高：曲线(2)}}.
\]
对同温不同气体，氢气分子质量小、速度尺度高，
因此$\boxed{\mathrm{H_2}:\text{曲线(2)},\quad
\mathrm{O_2}:\text{曲线(1)}}$。
\crosscheck 两条曲线都应归一化为1，峰高变化不能单独视为“分子总数变化”。

\sourceproblem{5.2.8}{125}{125}
\problemtext 一定量理想气体先等容加热至温度的2倍，
再等压膨胀至体积的2倍。求最终分子平均自由程与初始值之比
（温度变化不影响有效碰撞直径）。
\workings
分子数$N$不变，硬球模型下
$\overline\lambda=1/(\sqrt2\pi d^2n)$，
$n=N/V$。第一次体积不变，$\lambda$不变；
第二次体积变为$2V_0$，数密度减半，故
\[
\boxed{\overline\lambda_f/\overline\lambda_0=2}.
\]
\crosscheck 等压使二次温度再加倍，但自由程由数密度决定；不应把温度的总4倍直接当作自由程4倍。
% === END textbook-ch05-blanks.tex ===

% === BEGIN textbook-ch05-solutions-001-009.tex blob cc950d0555515df4b882113e32d5ace774362d10 ===
% 《大学物理简明教程》第5章解答题5.3.1—5.3.9，印刷/PDF第125—126页。
% 参数从两页扫描图片逐项核对，原创解析，待独立学术复核。
\subsection{解答题：分子动理论的定量应用}

\sourceproblem{5.3.1}{125}{125}
\problemtext 容器中的氧气压强$p=0.1\,\mathrm{MPa}$、温度$27^\circ\mathrm C$。
求分子数密度$n$、单个氧分子质量$m$和气体质量密度$\rho$。
\workings
采用$T=300.15$ K、$p=10^5$ Pa、$M_{\mathrm{O_2}}=0.028\,\mathrm{kg/mol}$：
\[
\boxed{n=\frac{p}{k_{\mathrm B}T}\approx
2.41\times10^{25}\,\mathrm{m^{-3}}},
\quad
\boxed{m=\frac{M_{\mathrm{O_2}}}{N_A}
\approx4.65\times10^{-26}\,\mathrm{kg}},
\]
\[
\boxed{\rho=nm=\frac{pM_{\mathrm{O_2}}}{RT}
\approx1.12\,\mathrm{kg/m^3}}.
\]
\crosscheck 将$n$乘单分子质量所得密度与理想气体状态方程求得的一致。

\sourceproblem{5.3.2}{125}{125}
\problemtext 一水银气压计原水银柱高$0.76\,\mathrm m$，玻璃管顶端至水银液面的空隙为$0.12\,\mathrm m$，
管横截面积$2.0\times10^{-4}\,\mathrm{m^2}$。注入少量He后水银柱高降至$0.60\,\mathrm m$。
温度$27^\circ\mathrm C$、水银密度$1.36\times10^4\,\mathrm{kg/m^3}$，求氦气质量（摩尔质量$0.004\,\mathrm{kg/mol}$）。
\workings
忽略原真空残余气体压强，注入后He的压力等于水银柱减少部分所对应的压差：
\[
p_{\mathrm{He}}=\rho_{\mathrm{Hg}}g(0.76-0.60)
=2.13248\times10^4\,\mathrm{Pa}.
\]
此时上方气体所占长度增加到
$0.12+(0.76-0.60)=0.28$ m，故
$V=(0.28)(2.0\times10^{-4})=5.6\times10^{-5}\,\mathrm{m^3}$。
\[
\boxed{m_{\mathrm{He}}
=\frac{p_{\mathrm{He}}VM_{\mathrm{He}}}{RT}
\approx1.91\times10^{-6}\,\mathrm{kg}}.
\]
\crosscheck 压差乘气体体积具有能量量纲；必须使用\emph{下降后的气柱高度}计算气体体积，不能沿用初始0.12 m。

\sourceproblem{5.3.3}{125}{125}
\problemtext 理想气体分子方均根速率为$400\,\mathrm{m/s}$，
压强为1 atm。求气体质量密度。
\workings
压强和方均根速率满足
$p=\frac13\rho\langle v^2\rangle$，故
\[
\boxed{\rho=\frac{3p}{v_{\mathrm{rms}}^2}
=\frac{3(101325)}{400^2}
\approx1.90\,\mathrm{kg/m^3}}.
\]
\crosscheck 压强若取$1.013\times10^5$ Pa，结果的有效位数不变。

\sourceproblem{5.3.4}{125--126}{125--126}
\problemtext 一瓶氧气和一瓶氢气处于相同温度；
已知氢气分子的平均平动动能
$\overline\varepsilon_{\mathrm H_2}=6.21\times10^{-21}\,\mathrm J$。
求氧气分子平均平动动能及方均根速率，并求氧气温度。
\workings
分子平均平动动能只由绝对温度决定，
故氧气分子也为$\boxed{6.21\times10^{-21}\,\mathrm J}$。
\[
T=\frac{2\overline\varepsilon}{3k_{\mathrm B}}
\approx\boxed{299.9\,\mathrm K}.
\]
氧气方均根速率：
\[
\boxed{v_{\mathrm{rms,O_2}}
=\sqrt{\frac{3RT}{0.032\,\mathrm{kg/mol}}}
\approx483\,\mathrm{m/s}}.
\]
\crosscheck 同温氢气的方均根速率约为氧气4倍，不能因两者平动平均动能相同就认为平均速率相同。

\sourceproblem{5.3.5}{126}{126}
\problemtext 外太空每立方厘米中约有一个氢原子，
温度$T=3.5\,\mathrm K$。求这些氢原子的方均根速率及其产生的压强。
\workings
数密度$n=1\,\mathrm{cm^{-3}}=10^6\,\mathrm{m^{-3}}$。
以氢原子质量$m_{\mathrm H}\approx1.66054\times10^{-27}\,\mathrm{kg}$，
\[
\boxed{v_{\mathrm{rms}}
=\sqrt{\frac{3k_{\mathrm B}T}{m_{\mathrm H}}}
\approx295\,\mathrm{m/s}},
\]
\[
\boxed{p=nk_{\mathrm B}T\approx4.83\times10^{-17}\,\mathrm{Pa}}.
\]
\crosscheck 这是低密度经典理想气体的形式估算；极稀薄宇宙空间的真实气体不必处于局部热平衡。

\sourceproblem{5.3.6}{126}{126}
\problemtext 体积$V=1\,\mathrm{m^3}$，其中含
$N_1=1.0\times10^{25}$个氢气分子和$N_2=4.0\times10^{25}$个氧气分子，
混合气体温度$400\,\mathrm K$。求分子的总平动动能及混合气体压强。
\workings
总分子数$N=N_1+N_2=5.0\times10^{25}$。由自由度均分定理与道尔顿定律：
\[
\boxed{E_{\mathrm{trans}}=\frac32Nk_{\mathrm B}T
\approx4.14\times10^5\,\mathrm J},
\quad
\boxed{p=\frac{Nk_{\mathrm B}T}{V}
\approx2.76\times10^5\,\mathrm{Pa}}.
\]
\crosscheck 即便两种气体分子质量不同，平均平动动能仍由共同温度确定；两种气体分压相加为总压。

\sourceproblem{5.3.7}{126}{126}
\problemtext 解释分子动理论中的六种能量表达式：
$\frac12k_{\mathrm B}T$、$\frac32k_{\mathrm B}T$、
$\frac i2k_{\mathrm B}T$、
$\frac{M}{M_{\mathrm{mol}}}\frac i2RT$、
$\frac i2RT$和$\frac32RT$的物理含义。
\workings
能量均分定理下，每个二次型自由度的平均能量为
$\boxed{\frac12k_{\mathrm B}T}$；
一个分子的三维平动动能均值为
$\boxed{\frac32k_{\mathrm B}T}$；
一个具有$i$个有效自由度的分子的平均总能量为
$\boxed{\frac i2k_{\mathrm B}T}$。
质量为$M$、摩尔质量$M_{\mathrm{mol}}$的理想气体有
$n=M/M_{\mathrm{mol}}$ mol，其内能为
$\boxed{U=n\frac i2RT}$。
$\frac i2RT$是1 mol同类气体的内能；
$\frac32RT$是1 mol任意经典理想气体的\emph{平动}总动能。
\crosscheck $k_{\mathrm B}T$是每分子尺度、$RT$是每摩尔尺度，$R=N_Ak_{\mathrm B}$，勿混用。

\sourceproblem{5.3.8}{126}{126}
\problemtext 容器内混合了11 kg二氧化碳和2 kg氢气，
题设将\emph{两种气体均近似为刚性双原子分子理想气体}；
已知混合气体内能$U=8.1\times10^6$ J。
求温度以及两种气体每分子的平均动能。
\workings
依\emph{题设有效自由度$i=5$}（尽管真实CO$_2$是线性三原子分子），
$n_1=11/0.044=250$ mol、$n_2=2/0.002=1000$ mol，
总摩尔数1250 mol。
\[
U=\frac52(n_1+n_2)RT
\quad\Longrightarrow\quad
\boxed{T=\frac{2U}{5(1250)R}\approx312\,\mathrm K}.
\]
两种分子的平均总动能相同，均为
\[
\boxed{\overline\varepsilon
=\frac52k_{\mathrm B}T\approx1.08\times10^{-20}\,\mathrm J}.
\]
\crosscheck 特别依赖题目要求的“五自由度”近似；真实CO$_2$的自由度激发程度依温度而异，此处不能擅自用真实分子结构替代教材题设。

\sourceproblem{5.3.9}{126}{126}
\problemtext 温度$27^\circ\mathrm C$时1 mol氢气分子的
平动动能、转动动能及内能分别是多少（刚性双原子模型）？
\workings
$T=300.15$ K，故
\[
\boxed{E_{\mathrm{trans}}=\frac32RT\approx3.74\times10^3\,\mathrm J},
\quad
\boxed{E_{\mathrm{rot}}=RT\approx2.50\times10^3\,\mathrm J},
\]
\[
\boxed{U=\frac52RT\approx6.24\times10^3\,\mathrm J}.
\]
\crosscheck 上述三者满足$E_{\mathrm{trans}}+E_{\mathrm{rot}}=U$。
% === END textbook-ch05-solutions-001-009.tex ===

% === BEGIN textbook-ch05-solutions-010-018.tex blob 2c51c7041c535069628f0069b02514bc370c6258 ===
% 《大学物理简明教程》第5章气体动理论，解答题5.3.10—5.3.18。
% 用户扫描PDF/书内页126，保留模型、数据与编号，独立LaTeX推导；待学术二校。
\subsection{解答题：能量、分布函数与碰撞自由程}

\sourceproblem{5.3.10}{126}{126}
\problemtext 一定量刚性双原子分子理想气体体积为
$V=2\times10^{-3}\,\mathrm{m^3}$，内能$U=6.75\times10^2$ J。
(1) 求压强；(2) 若总分子数$N=5.4\times10^{22}$，
求每分子平均平动动能和温度。
\workings
刚性双原子分子自由度$i=5$，由
$U=\tfrac52Nk_{\mathrm B}T$及$pV=Nk_{\mathrm B}T$，
\[
\boxed{p=\frac{2U}{5V}
=\frac{2(675)}{5(0.002)}
=1.35\times10^5\,\mathrm{Pa}}.
\]
平均平动动能
\[
\boxed{\overline\varepsilon_{\mathrm{trans}}
=\frac32k_{\mathrm B}T=\frac{3U}{5N}
=7.50\times10^{-21}\,\mathrm J}.
\]
温度
\[
\boxed{T=\frac{2U}{5Nk_{\mathrm B}}
\approx362\,\mathrm K}.
\]
\crosscheck $pV=(2/5)U$，两边均有焦耳量纲；
计算温度不需要知道气体分子质量。

\sourceproblem{5.3.11}{126}{126}
\problemtext 容积$1\,\mathrm{m^3}$的容器内有1 mol氢气（按刚性双原子分子）；
整个容器原以$10\,\mathrm{m/s}$平移，突然停下。
若氢气整体运动的机械动能中有80\%转化成分子热运动内能，求气体温度与压强的升高量。
\workings
1 mol氢气质量$m_{\mathrm H_2}\approx0.002\,\mathrm{kg}$，
整体平动动能
$\tfrac12mv^2=\tfrac12(0.002)(10^2)=0.10$ J。
进入气体内能的部分为
$\Delta U=0.8(0.10)=0.08$ J。
由$n=1$ mol与定容热容$C_V=\tfrac52R$，
\[
\boxed{\Delta T=\frac{\Delta U}{(5/2)R}
\approx3.85\times10^{-3}\,\mathrm K}.
\]
再利用$pV=nRT$，容积固定，
\[
\boxed{\Delta p=\frac{nR\Delta T}{V}
=\frac{2}{5}\frac{\Delta U}{V}
=0.032\,\mathrm{Pa}}.
\]
\crosscheck 如把80\%的转化误当100\%，压强会被高估至0.040 Pa；
气体升温虽很小，但不应为零。

\sourceproblem{5.3.12}{126}{126}
\problemtext 标准状态下氧气（刚性双原子分子）和氦气的体积之比
$V_{\mathrm{O_2}}/V_{\mathrm{He}}=1/2$，求两种气体内能之比。
\workings
标准状态两气体具有相同$T,p$，故$n_{\mathrm{O_2}}/n_{\mathrm{He}}
=V_{\mathrm{O_2}}/V_{\mathrm{He}}=1/2$。
氧气$i=5$，氦气$i=3$，
\[
\boxed{\frac{U_{\mathrm{O_2}}}{U_{\mathrm{He}}}
=\frac{(5/2)n_{\mathrm{O_2}}RT}
 {(3/2)n_{\mathrm{He}}RT}
=\frac{5}{6}}.
\]
\crosscheck 总内能比较既要计入分子自由度，又要计入不同容积导致的摩尔数差异。

\sourceproblem{5.3.13}{126}{126}
\problemtext 说明麦克斯韦速率分布函数$f(v)$的意义，并解释
$f(v)\mathrm dv$、$nf(v)\mathrm dv$、$Nf(v)\mathrm dv$、
$\int_0^v f(s)\mathrm ds$、$\int_0^\infty f(s)\mathrm ds$及
$\int_{v_1}^{v_2}Nf(v)\mathrm dv$。
其中$n$为分子数密度，$N$为整个系统的分子总数。
\workings
$f(v)$为每单位速率的\emph{概率密度}，单位为$\mathrm{s/m}$。
\[
\begin{array}{l|l}
\text{表达式}&\text{物理意义}\\\hline
f(v)\mathrm dv&\text{速率落在}[v,v+\mathrm dv]\text{的分子所占比例}\\
nf(v)\mathrm dv&\text{单位体积内该速率区间的分子数}\\
Nf(v)\mathrm dv&\text{整个系统内该速率区间的分子数}\\
\int_0^v f(s)\mathrm ds&\text{速率不大于}v\text{的分子比例}\\
\int_0^\infty f(s)\mathrm ds&1\text{（归一化）}\\
\int_{v_1}^{v_2}Nf(v)\mathrm dv&
\text{系统内速率在}[v_1,v_2]\text{的分子数}
\end{array}
\]
\crosscheck 只有概率、比例无量纲；分子数在统计模型中可按连续分布近似为期望数量。

\sourceproblem{5.3.14}{126}{126}
\problemtext 质量$m=6.2\times10^{-17}\,\mathrm{kg}$的微粒悬浮于
$27^\circ\mathrm C$的液体，观测到其方均根速率为$1.4\,\mathrm{cm/s}$。
按经典麦克斯韦分布，求阿伏伽德罗常量$N_A$。
\workings
先根据均方速度与能量均分定理求玻尔兹曼常量$k_{\mathrm B}$，
再通过已知普适气体常量$R=N_Ak_{\mathrm B}$求$N_A$。
令$v_{\mathrm{rms}}=0.014\,\mathrm{m/s}$、$T=300.15\,\mathrm K$，
由$\tfrac12m\langle v^2\rangle=\tfrac32k_{\mathrm B}T$得
\[
\boxed{k_{\mathrm B}=\frac{m v_{\mathrm{rms}}^2}{3T}
=\frac{(6.2\times10^{-17})(0.014)^2}{3(300.15)}
\approx1.35\times10^{-23}\,\mathrm{J/K}}.
\]
取$R=8.314462618\,\mathrm{J/(mol\,K)}$，
则
\[
\boxed{N_A=\frac{R}{k_{\mathrm B}}
\approx6.16\times10^{23}\,\mathrm{mol^{-1}}}.
\]
\crosscheck 本实验估计值与标准$6.022\times10^{23}\,\mathrm{mol^{-1}}$同数量级；
估算偏差来自给定实验数值舍入。原题需再核对其所称常量的中文用字。

\sourceproblem{5.3.15}{126}{126}
\problemtext 两瓶氧气、氢气处于相同压强与绝对温度，
氧气体积是氢气体积的两倍。
求两种气体的分子数密度之比和方均根速率之比。
\workings
理想气体分子数密度$n=p/(k_{\mathrm B}T)$，
故尽管两容器体积不同，
\[
\boxed{\frac{n_{\mathrm{O_2}}}{n_{\mathrm{H_2}}}=1}.
\]
对于方均根速率，
$v_{\mathrm{rms}}=\sqrt{3RT/M_{\mathrm{mol}}}$，故
\[
\boxed{\frac{v_{\mathrm{rms,O_2}}}{v_{\mathrm{rms,H_2}}}
=\sqrt{\frac{M_{\mathrm{H_2}}}{M_{\mathrm{O_2}}}}
=\sqrt{\frac{2}{32}}=\frac14}.
\]
\crosscheck 容积比2:1使\emph{总分子数}比为2:1，但\emph{数密度}仍为1:1，不得混淆。

\sourceproblem{5.3.16}{126}{126}
\problemtext $N$个粒子的速率分布图中，$Nf(v)$由原点线性升至
$(v_0,a)$，然后在$v_0\le v\le2v_0$恒等于$a$，区间外为零。
求$f(v)$、平均速率与速率位于$1.5v_0$到$2v_0$间的粒子数。
\workings
归一化要求图下面积为$N$：
\[
N=\frac12av_0+av_0=\frac32av_0
\Longrightarrow a=\frac{2N}{3v_0}.
\]
故
\[
\boxed{f(v)=
\begin{cases}
\frac{2v}{3v_0^2},&0\le v<v_0,\\
\frac{2}{3v_0},&v_0\le v\le2v_0,\\
0,&\text{其他}.
\end{cases}}
\]
平均速率
\[
\overline v=\int_0^\infty v f(v)\,dv
=\frac{2}{3v_0^2}\frac{v_0^3}{3}
+\frac{2}{3v_0}\frac{(2v_0)^2-v_0^2}{2}
=\boxed{\frac{11}{9}v_0}.
\]
所求区间的分子数
\[
\boxed{N_{[1.5v_0,2v_0]}
=\int_{1.5v_0}^{2v_0}Nf(v)\,dv
=\frac N3}.
\]
\crosscheck $\int f(v)\,dv=1$；区间概率为$1/3$不随$v_0$选取而变化。

\sourceproblem{5.3.17}{126}{126}
\problemtext 真空管内压强$p=1.38\times10^{-3}\,\mathrm{Pa}$，
温度$27^\circ\mathrm C$；气体分子有效直径$d=3.0\times10^{-10}\,\mathrm m$。
求分子数密度和平均自由程。
\workings
温度$T=300.15$ K，
\[
\boxed{n=\frac p{k_{\mathrm B}T}
\approx3.33\times10^{17}\,\mathrm{m^{-3}}}.
\]
硬球气体自由程
\[
\boxed{\overline\lambda
=\frac{1}{\sqrt2\pi d^2n}
=\frac{k_{\mathrm B}T}{\sqrt2\pi d^2p}
\approx7.51\,\mathrm m}.
\]
\crosscheck 当管道尺寸远小于计算自由程时，
实际粒子与器壁的碰撞可能远比粒子间碰撞频繁；
这里求的是硬球\emph{分子间}平均自由程。

\sourceproblem{5.3.18}{126}{126}
\problemtext 标准状态下氢气分子有效直径
$d=2.30\times10^{-10}\,\mathrm m$。
求平均自由程和平均碰撞频率。
\workings
取$p_0=101325$ Pa、$T_0=273.15$ K、
$M_{\mathrm{H_2}}=2.0\times10^{-3}\,\mathrm{kg/mol}$。
数密度
$n_0=p_0/(k_{\mathrm B}T_0)\approx2.687\times10^{25}\,\mathrm{m^{-3}}$。
\[
\boxed{\overline\lambda=
\frac{1}{\sqrt2\pi d^2n_0}
\approx1.58\times10^{-7}\,\mathrm m}.
\]
按经典麦克斯韦分布
$\overline v=\sqrt{8RT_0/(\pi M_{\mathrm{H_2}})}
\approx1.70\times10^3\,\mathrm{m/s}$，故
\[
\boxed{\overline Z=\frac{\overline v}{\overline\lambda}
\approx1.07\times10^{10}\,\mathrm{s^{-1}}}.
\]
\crosscheck $\overline v$必须用\emph{平均速率}而非方均根速率计算平均碰撞频率；不同的氢气摩尔质量舍入会引入约1\%量级的末位差异。
% === END textbook-ch05-solutions-010-018.tex ===
